\documentclass[11pt]{article}
\usepackage[margin=1in]{geometry}
\usepackage{amsmath,amssymb,amsthm}
\usepackage{graphicx}
\usepackage[colorlinks=true,linkcolor=blue,citecolor=blue,urlcolor=blue]{hyperref}

\newtheorem{theorem}{Theorem}
\newtheorem{lemma}{Lemma}

\newcommand{\R}{\mathbb{R}}
\newcommand{\cL}{\mathcal{L}}
\newcommand{\cN}{\mathcal{N}}
\newcommand{\Id}{\operatorname{Id}}
\newcommand{\tr}{\operatorname{tr}}
\newcommand{\im}{\operatorname{im}}
\newcommand{\kerop}{\operatorname{ker}}

\begin{document}

\begin{center}
{\Large A full answer to the odd-dimensional case in arXiv:1905.12257}\\[3mm]
{\large Candidate solution packet for \texttt{fa\_banach\_001}}\\[2mm]
\today
\end{center}

\section*{Source and Target}

The source paper is Kadets, Martin, Meri, Perez, and Quero,
\emph{On the numerical index with respect to an operator}, arXiv:1905.12257
\cite{KMMPQ}.  After Theorem 4.7 the authors prove the desired zero-value
result for $\cL(H)$ and $\mathcal K(H)$ when the real Hilbert space $H$ is
infinite-dimensional or even-dimensional, and then leave the finite odd case
open.

\begin{center}
\includegraphics[width=0.98\linewidth]{figures/open_problem_crop.png}
\end{center}

The question in context is whether, for a finite-dimensional real Hilbert
space $H$ of odd dimension at least $3$,
\[
  n_G(\cL(H),Y)=0
\]
holds for every Banach space $Y$ and every norm-one operator
$G\in \cL(\cL(H),Y)$.  Dimension $1$ is outside the context of Theorem 4.7 and
would be exceptional.

\section*{Idea of the Proof}

The even-dimensional proof in the source paper uses a surjective isometry
of $H$ with numerical radius zero, then transports it by left multiplication
to $\cL(H)$.  Such an invertible skew isometry cannot exist in odd real
dimension.  The replacement is to use two-sided infinitesimal isometries on
$\cL(H)$:
\[
  X\longmapsto AX+XB,
\]
where $A$ and $B$ are skew-adjoint on $H$.  These maps have numerical radius
zero because they generate the isometry group
$X\mapsto e^{tA}Xe^{tB}$.  By choosing the one-dimensional kernels of $A$ and
$B$ to match a single Hilbert-Schmidt orthogonality condition, every prescribed
operator $C\in\cL(H)$ lies in the range of one such zero-radius map.  This gives
exactly the range-covering hypothesis in Proposition 3.10(c) of the source
paper.

\section*{Theorem}

\begin{theorem}
Let $H$ be a real Hilbert space of odd dimension $n\geq 3$, and put
$W=\cL(H)$ with its operator norm.  Then
\[
  n_G(W,Y)=0
\]
for every Banach space $Y$ and every norm-one operator $G\in\cL(W,Y)$.
Consequently the question after Theorem 4.7 of \cite{KMMPQ} has an
affirmative answer.
\end{theorem}

\begin{proof}
By Proposition 3.10(c) of \cite{KMMPQ}, it is enough to prove that
\[
  \bigcup\{R(W): R\in\cL(W),\ v(R)=0\}
\]
is dense in $W$.  We prove the stronger statement that this union is all of
$W$.

Fix $C\in W$.  Choose a unit vector $b\in H$.  Since $n\geq 3$, we may choose
a unit vector $a\in H$ such that
\[
  \langle Cb,a\rangle=0;
\]
if $Cb=0$, choose any unit $a$.  Write $n=2m+1$.

Choose skew-adjoint operators $A,B\in\cL(H)$ such that
\[
  \kerop A=\R a,\qquad \kerop B=\R b,
\]
and such that the nonzero rotation frequencies of $A$ and $B$ are disjoint.
Concretely, decompose $H=\R a\oplus E_A$ and $H=\R b\oplus E_B$, split each
even-dimensional complement into orthogonal $2$-planes, let $A$ be block
rotation by positive frequencies $\alpha_1,\ldots,\alpha_m$ on $E_A$, let $B$
be block rotation by positive frequencies $\beta_1,\ldots,\beta_m$ on $E_B$,
and choose the two finite sets $\{\alpha_i\}$ and $\{\beta_j\}$ disjoint.

Define the Sylvester operator $R\in\cL(W)$ by
\[
  R(X)=AX+XB\qquad (X\in W).
\]
First, $v(R)=0$.  Indeed, because $A$ and $B$ are skew-adjoint,
$e^{tA}$ and $e^{tB}$ are orthogonal for every real $t$, and hence
\[
  U_t(X)=e^{tA}Xe^{tB}
\]
is a surjective isometry of $W$.  The left and right multiplication parts
commute, so $U_t=e^{tR}$.  The derivative formula for the numerical range
used in Lemma 2.4 of \cite{KMMPQ} gives
\[
 \sup \operatorname{Re} V(\cL(W),\Id_W,R)
 =\lim_{t\downarrow 0}\frac{\|\Id_W+tR\|-1}{t}.
\]
Since $e^{tR}$ is an isometry and $\Id_W+tR=e^{tR}+O(t^2)$ in operator norm,
the right-hand side is at most $0$.  Applying the same argument to $-R$ shows
that the real numerical range of $R$ is contained in $\{0\}$.  Therefore
$v(R)=0$.

It remains to show that $C\in R(W)$.  Use the Hilbert-Schmidt inner product
on the finite-dimensional real vector space $W$.  Because $A^*=-A$ and
$B^*=-B$, the Hilbert-Schmidt adjoint of $R$ is $R^*=-R$.  Thus
\[
  \im R=(\kerop R^*)^\perp=(\kerop R)^\perp .
\]
We now compute $\kerop R$.  After complexifying, the eigenvalues of $A$ are
$0,\pm i\alpha_1,\ldots,\pm i\alpha_m$, and the eigenvalues of $B$ are
$0,\pm i\beta_1,\ldots,\pm i\beta_m$.  On a rank-one map carrying a
$B$-eigenvector with eigenvalue $\mu$ to an $A$-eigenvector with eigenvalue
$\lambda$, the complexified Sylvester map has eigenvalue $\lambda+\mu$.
Since the nonzero frequency sets are disjoint, the only way to have
$\lambda+\mu=0$ is $\lambda=\mu=0$.  Hence
\[
  \kerop R=\R K,\qquad K(h)=\langle h,b\rangle a.
\]
The Hilbert-Schmidt orthogonality condition for $C$ is
\[
  \langle C,K\rangle_{\mathrm{HS}}=\tr(K^*C)=\langle Cb,a\rangle=0.
\]
Therefore $C\in(\kerop R)^\perp=\im R$.  Since $C\in W$ was arbitrary, every
element of $W$ lies in the range of some zero-numerical-radius operator on
$W$.  Proposition 3.10(c) of \cite{KMMPQ} now yields
$n_G(W,Y)=0$ for every Banach space $Y$ and every norm-one
$G\in\cL(W,Y)$.
\end{proof}

\section*{Verification Notes}

The proof above is exact finite-dimensional linear algebra.  The included
script \path{code/verify_sylvester_cover.py} only sanity-checks the
Sylvester range computation numerically.  The command
\begin{verbatim}
cd runs/fa_banach_001/solutions/full/\
  1905.12257_odd_dim_operator_space_zero_index
conda run --no-capture-output -n sandbox python \
  code/verify_sylvester_cover.py --dims 3 5 7 --trials 20
\end{verbatim}
returned ranks $8,24,48$, i.e. $n^2-1$, with residuals below $1.2\cdot 10^{-13}$.
This is consistent with the proof that the only obstruction is the line
spanned by $a\otimes b^*$, but the computation is not used as evidence in
place of the proof.

\section*{Novelty and Scope}

Bounded local and web checks were performed on 2026-06-26.  The local run
indexes were searched for arXiv:1905.12257, ``numerical index with respect'',
``odd dimension'', and the operator-space phrasing.  Web/arXiv searches used
the exact phrases ``When H has odd dimension'', ``numerical index with respect
to an operator'' with ``odd dimension'' and ``L(H)'', the source title with
``odd dimension'', and the displayed expression
$n_G(\mathcal L(H),Y)=0$.  No later paper or existing packet answering this
exact question was found; the only direct web hit was the source arXiv record.

The packet answers only the odd-dimensional Hilbert-operator-space question
after Theorem 4.7.  The other questions in the same source paper, such as the
countability questions for $\mathcal N(Z)$, are not addressed here.

\begin{thebibliography}{9}
\bibitem{KMMPQ}
V. Kadets, M. Martin, J. Meri, A. Perez, and A. Quero,
\emph{On the numerical index with respect to an operator},
arXiv:1905.12257, 2019.
\end{thebibliography}

\end{document}
